\section{Discussion}
\label{sec:discussion}

\subsection{Geometry calibration and model discrepancy}
The length fit constrains one combination of source strength, material continuation and observation mapping. It leaves the B fusion envelope 19.54~\um{} wider and 4.88~\um{} shallower than the nominal references, and the fixed-parameter B-to-C shape response differs from the experimental means (Fig.~\ref{fig:comparison}). These discrepancies show which information the fit leaves unresolved: penetration, transverse heat distribution and the temperature field are not independently constrained by the surface length. Source representation, transport closure and the distinction between radiance-derived and temperature-threshold geometry remain possible contributors.

Geometric agreement can conceal thermal differences even in a more detailed model. In the bare-plate 316L calculations of Myers et al.~\cite{myers2023}, several Fresnel/accommodation parameter combinations gave similar agreement with cross-sectional area and width-to-depth ratio but different surface temperature profiles. Two-color measurements narrowed those combinations. The present comparison has a different design: only the LL length is fitted, and the other continuations retain its source factor. Their thermal differences measure response after one calibration; they do not demonstrate thermal differences between equally refitted branches.

The closest AMMT comparison also makes source and transport choices consequential. Kollmannsberger et al.~\cite{kollmannsberger2019} used measured laser input and an effective directional conductivity calibrated at B to improve the A/C geometric comparisons, while explicitly limiting the validity of internal melt-pool temperatures. Their result indicates a route to better geometry under a different source and closure, rather than a ranking against this isotropic Gaussian calculation. The current continuation experiment is useful precisely because it preserves those choices and asks what changes when the missing high-temperature property range is treated differently.

\subsection{Conductivity and storage in the thermal tail}
Keeping the calibrated source fixed reveals opposing property responses that a joint replacement would conceal. Conductivity holding extends the pool at either storage setting, while sensible-storage holding shortens it at either conductivity setting (Table~\ref{tab:factorial}). Their combination produces a 34.28~\um{} length increase, with the interaction modifying the magnitudes. The rear profiles locate the dominant response: the rear solidus accounts for 98.3\% of the conductivity-induced length increase, and the liquidus moves almost the same distance. The principal change is a displacement of the thermal tail relative to the source, rather than comparable widening of the rear phase interval (Fig.~\ref{fig:factorial}). None of the four corners removes the B wide/shallow discrepancy.

The adopted enthalpy law helps explain why sensible-storage holding has the smaller response. Inside the phase interval, the fixed latent term contributes 4666.67~\si{\joule\per\kilogram\per\kelvin} to effective specific heat. At 1320~$^\circ$C, holding sensible $c_p$ reduces it by 7.09\%, but reduces latent-inclusive effective specific heat by only 0.95\% and enthalpy from the cold reference by 0.66\%. The conductivity reduction at the same temperature is 22.48\%. These comparisons show that the two nominal property changes perturb transport and total storage by different amounts. They do not isolate latent heat as the cause: latent heat is unchanged in all corners, and each continuation starts below the solidus, so pre-melting enthalpy and the temperature--enthalpy inverse also change.

The diagnostics support a changed balance between coordinate enthalpy transport and diffusion, while ruling out a simple bulk-flux explanation. The selected axial-face P\'eclet median rises from 4.44 to 5.77 under conductivity holding. Yet outward diffusive power from the rear hot region rises from 14.724 to 14.913~W. Lower conductivity therefore does not imply lower integrated outward heat flow in this comparison; both the temperature gradients and the boundary of the temperature-selected region change. Surface recovery also changes consistently through $K(T)$ and its inverse. The resolved tail displacement is a response of the complete conduction model and observation operator. Assigning it solely to suppressed bulk heat escape would require a matched-region flux comparison or a separation of field and recovery contributions.

\subsection{Spatial phase span and material cooling time}
Rear-boundary location, phase separation and material time describe different features of a fitted field. The current-mesh B/C rear phase spans are 80.26 and 80.79~\um{}, while their passage times are 100.33 and 67.33~\us{} because scan speed rises from 0.8 to 1.2~m/s (Fig.~\ref{fig:trends}). Thus nearly unchanged spatial separation coexists with a much shorter material interval. The 49.02\% increase in interval-mean cooling rate follows from the same thresholds and the division by speed in Eq.~\eqref{eq:passage}; it is not an additional independent thermal test. Likewise, translating the two rear boundaries together can change total pool length substantially without producing a similar change in passage time.

Thermography distinguishes these questions experimentally. Hooper~\cite{hooper2018} measured both centerline spatial profiles and temperature histories at fixed build-plate locations, including pulsed scans and turn events. Those observations show why a spatial profile and a thermal history cannot be exchanged without accounting for the scan motion. Lane's cooling-rate exemplars~\cite{lane2020} use a below-solidus interval and a different rear-boundary definition, so they do not validate the present 1350--1290~$^\circ$C interval. A comparison would require a temperature trajectory over the same interval, with scan position, exposure and observation mapping included.

The passage transformation describes stationary material crossing the assumed steady field. It does not track liquid motion through that field or specify a nonequilibrium freezing path. This boundary matters when thermal outputs are used for microstructure: Plotkowski et al.~\cite{Plotkowski2017Rapid} treated melt-pool geometry verification separately from solidification-condition calculations and an experimental grain-structure comparison. Their AlSi10Mg/IN718 study provides a useful methodological distinction, not validation of this IN625 model. Here, neither grain morphology nor segregation has been tested, and the rear passage interval remains a model diagnostic.

\subsection{Thermal interpretation and discriminating evidence}
The shape residual does not identify a unique missing mechanism. An alternative source is a substantive explanation: Coleman et al.~\cite{dynamic2024} calibrated a dynamic volumetric source for bare IN625 at 195~W and 800~mm/s over several spot sizes, then applied it to a layer case. Their source adjusts depth and effective absorption with the local melt-pool state; their phase-dependent properties and 1410~K eutectic endpoint also differ from the present equilibrium interval. A comparison using the same measured beam, phase assumptions, calibration target and geometry operators would distinguish source redistribution from the effect of a different material closure. An imaging replay is a cheaper first test of how the radiance-to-threshold mismatch changes the fitted source factor. Neither comparison is supplied by the present continuation corners.

Two physical limits govern the thermal interpretation. First, the supplier property tables end at 982~$^\circ$C for $k$ and 1093~$^\circ$C for $c_p$, below the solidus~\cite{specialmetals625}. The L/H branches are continuation scenarios, not measured liquid-state functions. Second, the model excludes evaporation, melt flow and surface evolution. Its LL, HL, LH and HH surface peaks are approximately 2697, 3507, 2756 and 3721~$^\circ$C. In particular, the held-conductivity peaks cannot be interpreted as credible physical-temperature predictions from this balance. The model nevertheless exposes how those assumptions affect the fitted geometry and rear thresholds.

Numerical checks support the larger responses without removing these physical limits. The main 34--43~\um{} length changes exceed the inspected baseline axial sensitivity, while the submicrometre interactions lack branch-specific three-dimensional refinement evidence. The analytical comparison, energy balance and solver correction test numerical execution~\cite{oberkampf2002}; they do not validate liquid properties or a thermal trajectory. The experimental comparisons are also retrospective and combine length and cross-section quantities from separate observation populations, with descriptive uncertainty scales rather than a joint likelihood.

A discriminating thermal test would retain a defined source and observation mapping, reserve tracks or profiles before fitting, and compare an excluded temperature trajectory together with width and depth. Composition- and state-matched high-temperature property evidence would address which continuation is physically appropriate; a matched thermal trajectory would address whether the calibrated field is credible. Until those tests are available, the result is specific but useful: one geometric fit leaves distinct questions about fusion shape, tail position, phase separation and material cooling time, and the fixed-source continuation contrast reveals which of those questions the assumed properties influence.
